\documentclass[11pt,a4paper]{article}

% --------------------------------------------------
% Pacotes Essenciais (Padrão UFPR)
% --------------------------------------------------
\usepackage[utf8]{inputenc}
\usepackage[T1]{fontenc}
\usepackage[brazil]{babel}
\usepackage{geometry}
\usepackage{setspace}
\usepackage{graphicx}
\usepackage{helvet} % Fonte similar a Arial
\usepackage{indentfirst} % Indenta o primeiro parágrafo
\usepackage{url}

% Configuração de Fonte (Arial)
\renewcommand{\familydefault}{\sfdefault}

% Configuração de Margens (2cm em tudo)
\geometry{top=2cm, bottom=2cm, left=2cm, right=2cm}

% Espaçamento 1.5
\onehalfspacing

% Remove numeração de páginas automática (pode ajustar se necessário)
\pagenumbering{gobble}

\begin{document}

% --------------------------------------------------
% 1. CAPA
% --------------------------------------------------
\thispagestyle{empty}

\begin{center}
    \textbf{UFPR} \\
    \textbf{UNIVERSIDADE FEDERAL DO PARANÁ} \\
    \vspace{0.5cm}
    MINISTÉRIO DA EDUCAÇÃO \\
    UNIVERSIDADE FEDERAL DO PARANÁ \\
    PRÓ-REITORIA DE PESQUISA E PÓS-GRADUAÇÃO \\
    COORDENADORIA DE INICIAÇÃO CIENTÍFICA E INTEGRAÇÃO ACADÊMICA \\
    PROGRAMA DE INICIAÇÃO CIENTÍFICA E EM DESENVOLVIMENTO TECNOLÓGICO E INOVAÇÃO
    
    \vspace{3cm}
    
    \textbf{ANGELO MAN DE OLIVEIRA}
    
    \vspace{3cm}
    
    \textbf{RELATÓRIO PARCIAL}
    
    \vspace{1.5cm}
\end{center}

% Bloco de seleção (PIBIC / Modalidade)
\noindent
\begin{minipage}[t]{0.48\textwidth}
    \textbf{PROGRAMA DE IC:}\\
    (X) PIBIC\\
    ( ) PIBIC Af\\
    ( ) PIBIC EM\\
    ( ) PIBITI
\end{minipage}%
\begin{minipage}[t]{0.48\textwidth}
    \textbf{MODALIDADE:}\\
    ( ) CNPq\\
    ( ) UFPR TN\\
    ( ) Fundação Araucária\\
    ( ) Voluntária
\end{minipage}

\vspace{2cm}

% Bloco de Informações do Projeto (Canto Inferior Direito)
\begin{flushright}
\begin{minipage}{0.6\textwidth}
    \small
    \singlespacing
    Relatório apresentado à Coordenação de Iniciação Científica e Tecnológica da Universidade Federal do Paraná como requisito parcial da conclusão das atividades de Iniciação Científica – Edital 2025/2026.
    
    \vspace{0.3cm}
    \textbf{Orientador(a):} Prof.(a). Rodrigo Clemente Thom de Souza
    
    \vspace{0.3cm}
    \textbf{Título do Projeto:} Análise do alinhamento entre Reinforcement Pre-Training e Reinforcement Learning with Internal Feedback em tarefas de matemática e programação.
\end{minipage}
\end{flushright}

\vfill

\begin{center}
    \textbf{JANDAIA DO SUL} \\
    \textbf{2026}
\end{center}

\newpage

% --------------------------------------------------
% 2. TÍTULO (Página 2)
% --------------------------------------------------
\begin{center}
    \textbf{\large Análise do alinhamento entre Reinforcement Pre-Training e Reinforcement Learning with Internal Feedback em tarefas de matemática e programação}
\end{center}
\vspace{1cm}

% --------------------------------------------------
% 3. ALTERAÇÕES REALIZADAS (Obrigatório)
% --------------------------------------------------
\section{ALTERAÇÕES REALIZADAS NO PLANO DE TRABALHO SUBMETIDO}

Inicialmente, era previsto o uso dos datasets OmniMATH e Codeforces Problems, com o objetivo de abranger problemas matemáticos e algorítmicos. No entanto, a aplicação desses datasets implicaria em dificuldades práticas de avaliação automática.

Diante disso, optou-se por empregar o benchmark HumanEval para a avaliação do desempenho algorítmico e o benchmark GSM8K para a avaliação do desempenho matemático.
% --------------------------------------------------
% 4. RESUMO
% --------------------------------------------------
\section{RESUMO}

Este projeto tem como objetivo analisar o alinhamento entre Reinforcement Pre-Training (RPT) e Reinforcement Learning from Internal Feedback (RLIF) em tarefas de raciocínio matemático e programação. Para isso, são utilizados problemas dos datasets GSM8K e Humaneval, avaliando modelos de linguagem de pequeno porte sob diferentes regimes de aprendizado por reforço. As métricas de desempenho incluem acurácia em matemática, porcentagem de códigos aceitos (métrica Pass@1), estabilidade do treinamento e custo computacional. Até o momento, foram definidos os datasets principais, as métricas de avaliação, os modelos a serem comparados e desenvolvidos experimentos preliminares.


% --------------------------------------------------
% 5. INTRODUÇÃO
% --------------------------------------------------
\section{INTRODUÇÃO}

Os modelos de linguagem têm apresentado avanços significativos na resolução de problemas algorítmicos e matemáticos que exigem raciocínio complexo. Técnicas de treinamento como o pre-training e o aprendizado por reforço aplicado no fine-tuning têm se consolidado como ferramentas robustas para o aprimoramento do desempenho desses modelos.

Entretanto, as abordagens tradicionais do reinforcement learning, como o RLHF e o RLVR, no fine-tuning apresentam limitações relevantes. O RLHF depende de extensas anotações humanas, o que eleva custos e pode introduzir vieses, enquanto o RLVR exige verificadores específicos de domínio e respostas de referência, restringindo sua aplicação a cenários cuidadosamente curados, como matemática e programação.

Além disso, o pre-training tradicional apresenta limitações quando aplicado a tarefas que exigem raciocínio estruturado e tomada de decisão. Por ser baseado predominantemente na predição estatística do próximo token, esse processo não incentiva explicitamente a construção de cadeias de raciocínio nem a validação de hipóteses intermediárias.

Diante disso, este trabalho investiga o alinhamento entre RPT e RLIF em tarefas de matemática e programação, analisando o potêncial desse alinhamento.


% --------------------------------------------------
% 6. REVISÃO DA LITERATURA
% --------------------------------------------------
\section{REVISÃO DE LITERATURA}

O treinamento de modelos de linguagem por meio de aprendizado por reforço tem avançado rapidamente, gerando diferentes estratégias para superar limitações do ajuste supervisionado. Entre as abordagens emergentes destacam-se o \textit{Reinforcement Pre-Training} (RPT) e o \textit{Reinforcement Learning from Internal Feedback} (RLIF), que diferem essencialmente na origem e na natureza do sinal de recompensa utilizado.

O RPT propõe a incorporação de objetivos de reforço já na etapa de pré-treinamento, permitindo que o modelo adquira políticas e vieses inductivos úteis antes do fine-tuning em tarefas específicas \cite{dong2025rpt}. Essa estratégia pode favorecer a emergência de capacidades de raciocínio estruturado e reduzir a necessidade de grandes quantidades de supervisão posterior.

No cenário de ajuste fino, abordagens baseadas em recompensas verificáveis externas — que chamamos aqui de RLVR (Reinforcement Learning with Verifiable Rewards) — utilizam sinais objetivos como execução de código, testes unitários ou rótulos humanos verificáveis para orientar o aprendizado. Tais sinais tendem a produzir melhorias robustas em tarefas cuja correção pode ser automaticamente avaliada, especialmente em programação e raciocínio simbólico, mas dependem da disponibilidade e qualidade dos validadores externos.

Em contraste, o RLIF substitui ou complementa recompensas externas por sinais internos gerados pelo próprio modelo (por exemplo, medidas de confiança, consistência entre gerações ou estatísticas da distribuição preditiva) \cite{zhao2025rlif}. Métodos que exploram esse tipo de feedback intrínseco (como propostas baseadas em self-likelihood ou self-certainty) visam tornar o treinamento mais autônomo e escalável em cenários onde avaliações externas são inexistentes, caras ou ruidosas. No entanto, sinais internos podem introduzir vieses e requerem mecanismos de regularização para evitar colapsos de entropia ou deriva indesejada da política.

Paralelamente, avanços em algoritmos de otimização para RL em LLMs têm mostrado impacto prático. O Group Relative Policy Optimization (GRPO), apresentado no contexto do DeepSeekMath, é um exemplo de variante que compara respostas dentro de grupos de amostras para melhorar a estabilidade e a eficiência computacional em tarefas de raciocínio matemático \cite{shao2024deepseekmath}. A combinação de técnicas de otimização (p.ex., GRPO) com esquemas de recompensa apropriados tem sido crucial para resultados de ponta em problemas complexos sem recorrer a processos de votação massiva ou tool-use.

Em suma, a literatura aponta para uma transição parcial de metodologias fortemente dependentes de sinais externos para esquemas mais autônomos que exploram feedback interno, mas ressalta a necessidade de estudos comparativos que avaliem trade-offs entre desempenho, estabilidade e preservação do conhecimento prévio. É nesse contexto que o presente trabalho se insere: comparar empiricamente RPT, RLVR e RLIF em benchmarks de raciocínio matemático e geração de código, avaliando não apenas acurácia mas também estabilidade do treinamento e custo computacional.

% --------------------------------------------------
% 7. MATERIAS E METODOS
% --------------------------------------------------
\section{Materiais e Métodos}

\subsection{Modelo base}
Foi utilizado como modelo inicial o \textbf{Qwen2.5-Coder-1.5B-Instruct}, um modelo de linguagem causal especializado em geração de código. O ajuste fino foi realizado por meio de \textit{Low-Rank Adaptation} (LoRA), mantendo os pesos principais congelados para reduzir custo computacional.

\subsection{Paradigma de treinamento}
O modelo foi ajustado utilizando \textbf{Reinforcement Learning from Internal Feedback (RLIF)}, implementado com \textit{Group Relative Policy Optimization} (GRPO).

Diferentemente de abordagens supervisionadas ou baseadas em verificação externa, a função de recompensa foi definida exclusivamente a partir de um sinal interno do próprio modelo, baseado na probabilidade atribuída aos tokens gerados.

A recompensa foi definida como a média da log-probabilidade dos tokens da resposta gerada (self-likelihood), excluindo os tokens do prompt. Esse mecanismo caracteriza RLIF por utilizar feedback intrínseco derivado da distribuição preditiva do próprio modelo.

\subsection{Dados de treinamento}
Foi utilizado o dataset \textbf{MBPP (Mostly Basic Python Problems), versão sanitized}, contendo problemas de programação em linguagem natural acompanhados de implementações de referência em Python.

Os exemplos foram formatados em estrutura ChatML, instruindo o modelo a gerar exclusivamente o corpo indentado da função correspondente ao problema.

\subsection{Configuração de treinamento}

\begin{itemize}
\item Learning rate: $1\times10^{-6}$
\item LoRA rank: 16
\item Batch efetivo: 8 (acumulação de gradiente)
\item Gerações por prompt: 4
\item Penalidade KL: 0.01
\item Precisão FP16
\item Temperatura de geração: 0.7
\end{itemize}

O treinamento foi realizado por uma época completa sobre o conjunto MBPP.

\subsection{Avaliação}
O desempenho foi avaliado no benchmark \textbf{HumanEval}, utilizando a métrica \textbf{pass@1}, que mede a proporção de problemas resolvidos corretamente em uma única geração.

Foram comparados três modelos:

\begin{enumerate}
\item modelo base sem ajuste adicional
\item modelo ajustado com RLVR (baseline supervisionado)
\item modelo ajustado com RLIF (proposto)
\end{enumerate}

\section{Resultados e Discussão (Preliminares)}

\subsection{Desempenho global}

\begin{center}
\begin{tabular}{lc}
\hline
Modelo & Pass@1 \\
\hline
RLIF & \textbf{62.8\%} \\
Base & 61.0\% \\
RLVR & 60.4\% \\
\hline
\end{tabular}
\end{center}

O modelo treinado com RLIF apresentou desempenho ligeiramente superior ao modelo base e ao modelo ajustado supervisionadamente. A diferença observada corresponde a um ganho absoluto de aproximadamente 1--2 pontos percentuais.

\subsection{Interpretação do resultado}

O desempenho do modelo treinado com RLIF permaneceu próximo ao do modelo base, com uma melhora moderada observada nesta execução específica do experimento.

Como a avaliação foi realizada com uma única geração por problema e uma única execução de treinamento, não é possível afirmar significância estatística do ganho observado. Portanto, os resultados devem ser interpretados como evidência preliminar de que o uso de feedback interno não degradou o desempenho e pode produzir melhora marginal.

\subsection{Estabilidade relativa do comportamento}

A proximidade entre os resultados do modelo RLIF e do modelo base sugere que o treinamento com recompensa interna preservou, em grande medida, o comportamento funcional previamente aprendido durante o pré-treinamento e ajuste instrucional.

Isso indica que o sinal de recompensa interno utilizado no treinamento não induziu mudanças drásticas na política de geração sob o regime experimental adotado.

\subsection{Comparação com RLVR}

O modelo RLVR apresentou desempenho ligeiramente inferior ao modelo base nesta configuração experimental. Isso sugere que, para o conjunto de hiperparâmetros e dados utilizados, o treinamento supervisionado adicional não produziu ganho mensurável em pass@1.

Por outro lado, o modelo RLIF manteve desempenho comparável ao modelo base, apresentando pequena vantagem numérica.

\subsection{Limitações experimentais}

Os resultados ainda são preliminares devido a:

\begin{itemize}
\item avaliação com uma única amostra por problema
\item ausência de múltiplas execuções com diferentes sementes
\item ausência de intervalos de confiança estatísticos
\item treinamento realizado por apenas uma época
\end{itemize}

Consequentemente, o ganho observado deve ser interpretado como indicativo, e não conclusivo.

\subsection{Síntese preliminar}

Os experimentos indicam que:

\begin{enumerate}
\item o treinamento com RLIF preserva desempenho próximo ao modelo base
\item não houve evidência de degradação funcional sob recompensa interna
\item foi observada pequena melhora em avaliação de execução única
\item são necessários experimentos adicionais para confirmar robustez do efeito
\end{enumerate}


% --------------------------------------------------
% 9. REFERÊNCIAS
% --------------------------------------------------
\section{REFERÊNCIAS}

\begin{thebibliography}{99}

\bibitem{zhao2025rlif}
ZHAO, Xuandong; KANG, Zhewei; FENG, Aosong; LEVINE, Sergey; SONG, Dawn.
\textit{Learning to Reason without External Rewards}.
arXiv preprint arXiv:2505.19590, 2025.
Disponível em: \url{https://arxiv.org/abs/2505.19590}. 

\bibitem{dong2025rpt}
DONG, Qingxiu; DONG, Li; TANG, Yao; YE, Tianzhu; SUN, Yutao; SUI, Zhifang; WEI, Furu.
\textit{Reinforcement Pre-Training}.
arXiv preprint arXiv:2506.08007, 2025.
Disponível em: \url{https://arxiv.org/abs/2506.08007}. 

\bibitem{shao2024deepseekmath}
SHAO, Zhihong; WANG, Peiyi; ZHU, Qihao; XU, Runxin; SONG, Junxiao; BI, Xiao; 
ZHANG, Haowei; ZHANG, Mingchuan; LI, Y. K.; WU, Y.; GUO, Daya.
\textit{DeepSeekMath: Pushing the Limits of Mathematical Reasoning in Open Language Models}.
arXiv preprint arXiv:2402.03300, 2024.
Disponível em: \url{https://arxiv.org/abs/2402.03300}. 

\end{thebibliography}

\noindent

\end{document}
